% Draft for Springer's journal Virtual Reality. Rules: wiki/src/XR-training.md -> Journal rules.
% Every \tbd{} is a value or decision that does not exist yet. Results come from paper/xr/results.yaml only.
% Build: cd paper/xr && latexmk   (PDF in paper/xr/build/). Move into the Springer Nature LaTeX template before submitting.
\documentclass[11pt]{article}
\usepackage[T1]{fontenc}
\usepackage[utf8]{inputenc}
\usepackage[american]{babel}
\usepackage{csquotes}
\usepackage[style=authoryear,backend=biber,maxcitenames=2,uniquename=false]{biblatex}
\addbibresource{../references.bib}
\usepackage{amsmath,booktabs,tabularx,xcolor,geometry}
\usepackage[hidelinks]{hyperref}
\geometry{margin=25mm}
\newcommand{\tbd}[1]{\textcolor{red}{\textbf{[TBD: #1]}}}

\title{One Web Page, Two Headsets: Augmented and Virtual Reality Training for an Open-Source Gonioreflectometer}
\author{Aleksander Norman \and \tbd{co-authors and order}\\[2pt]
\small \tbd{affiliations: department, institution, city, state, country}\\
\small Corresponding author: \tbd{name, e-mail, ORCID}}
\date{Draft of 2026-11-25. No headset measurement has been made.}

\begin{document}
\maketitle

\begin{abstract}
A gonioreflectometer measures how a surface reflects light in every direction, and students find the instrument hard
to picture: four rotary axes, travel limits that protect the camera, and a measured quantity defined by two angles
that are not the motor angles. We present a training tool that draws these ideas in place. It overlays the travel of
each axis and its limits on the real machine, traces the light from the lamp to the sample and on to the camera, and
shows the reflectance lobe that the measurement samples. The tool is a single web page built on the WebXR Device API.
The same code and content run in augmented reality on a Microsoft HoloLens 2 and a Meta Quest, in virtual reality
around a model of the machine, and in a desktop browser. Geometry and the reflectance model are taken from the
instrument's own design files and analysis software, and automated tests hold the two together. Because both headsets
run identical content, differences we measure belong to the platforms. We define a comparison protocol covering frame
timing, registration error, drift, field of view, legibility, learning gain, usability, workload and discomfort.
\tbd{two sentences of measured results; none exist yet}. The tool and the instrument are open source.
\end{abstract}

\noindent\textbf{Keywords:} augmented reality; virtual reality; WebXR; training; gonioreflectometer; reflectance

\section{Introduction}
The bidirectional reflectance distribution function (BRDF) describes how much light a surface sends toward a viewer
for each direction of illumination \parencite{G1}. A gonioreflectometer measures it by moving a light source, a sample
and a detector through many angle pairs \parencite{R1}. OpenGonio is an open-source, printable instrument of this kind with four
motorised axes: a light arm (M4), a table (M3), a tilt (M2) and a spin (M1).

New users make the same mistakes. They confuse motor angles with the angles of incidence and reflection, they do not
know where an axis will stop, and they expect the camera to see a highlight where there is none. A diagram shows one
pose. The machine itself shows no light path. Augmented reality (AR) can place the missing information on the
instrument, and virtual reality (VR) can offer the same lesson to a student who has no instrument
\parencite{X8,X12,X13}.

This paper makes three contributions:
\begin{enumerate}
\item a training tool that overlays axis travel, light paths and the BRDF lobe on a real gonioreflectometer, generated
 from the instrument's design files;
\item an implementation as one web page that runs unchanged on an optical see-through and a video see-through headset;
\item a protocol, and \tbd{its measured results}, comparing the two platforms on the same content.
\end{enumerate}

\section{Related work}
\subsection{Mixed reality displays and registration}
We use the terms of \textcite{X7}. \textcite{X8} separates optical see-through displays, which add light to the user's
view, from video see-through displays, which show a camera image, and names registration as the central problem of
AR. The two headsets compared here are one of each kind.

\subsection{Augmented reality for learning}
Reviews report learning benefits of AR for spatial content and note usability problems and attention costs
\parencite{X12,X13}. \tbd{closest prior work on AR or VR training for optical instruments; search and cite before
submission}.

\subsection{Web delivery}
The WebXR Device API gives a web page stereo views, tracking and input on a headset \parencite{X1}; its AR module adds
sessions blended with the real environment \parencite{X2}. Both vendors document support
\parencite{X3,X4}. The HoloLens 2 is out of production with updates until the end of 2027 \parencite{X5}, and users
have reported periods in which its browser did not start AR sessions \parencite{X6}. We treat this as a risk to the
study, not as a result.

\section{System design}
\subsection{Instrument and coordinate frames}
The machine frame has its origin on the deck, $z$ up, and a zero point on the sample face 150~mm above it. The camera
is fixed on the $+x$ axis, 350~mm from the zero point. The lamp rides the light arm on a circle of radius 300~mm at
azimuth $M_4$, limited to $20^\circ$--$324^\circ$ so that it cannot reach the camera. The sample orientation is
\begin{equation}
R = R_z(M_3)\,R_y(-M_2)\,R_x(M_1), \qquad \mathbf{n} = R\,\mathbf{e}_x ,
\end{equation}
and the angles of incidence and reflection are $\theta_i = \arccos(\mathbf{n}\cdot\mathbf{l})$ and
$\theta_r = \arccos(\mathbf{n}\cdot\mathbf{v})$ for unit vectors $\mathbf{l}$ to the lamp and $\mathbf{v}$ to the camera.

\subsection{Lessons}
Five lessons each add one idea: (1) the travel of every axis and its limits; (2) the incident ray and $\theta_i$;
(3) the mirror direction, the camera ray and $\theta_r$; (4) the BRDF lobe for the current lamp position; (5) a sweep
of the lamp in which one bar per lamp angle shows the camera reading. Lessons animate without input; one pinch or
trigger press advances to the next.

\subsection{Reflectance model}
The lobe is the model the instrument's software fits to measured data: a Lambertian term plus a GGX microfacet term
\parencite{R6} with base colour, roughness and metallic parameters. Its radius in direction $\boldsymbol\omega$ is
$\sqrt{f(\mathbf{l},\boldsymbol\omega)\cos\theta_r}$, a square-root scale that keeps narrow and wide lobes visible.

\subsection{Registration}
In AR the user pinches at the centre of the sample and then at the camera lens. The first point fixes the zero point,
the horizontal direction to the second fixes the heading, and gravity fixes the vertical. The difference between the
distance of the two pinches and the designed 350~mm is logged as a residual. The method needs no marker and no access
to the headset's cameras. It does not correct tilt.

\subsection{Logging}
Each session records the time between frames, the number of views, the framebuffer size, the reference space and the
alignment residual, and reports mean and median frame time, 95th and 99th percentiles, and the count of frames longer
than 1.5 times the median.

\section{Implementation}
The tool is one HTML page and one script of about 300 lines, drawn with WebGL and no engine. Each frame it rebuilds
the overlay as triangles and draws it once per eye with the view and projection matrices supplied by the headset. Lines
are two crossed ribbons with a width in millimetres, blended with premultiplied alpha so that the same colours are
correct on an additive display and on a video display. In VR and in the desktop view the CAD model of the machine is
drawn and posed with the same axis rotations as the design's interference check; in AR it is hidden.

Axis positions, travel limits, and the lamp and camera points are read from a file generated by the CAD build. A
project check fails if that file disagrees with the machine configuration. Automated tests run the page's own
functions and compare the reflectance model with the analysis software to a relative tolerance of $10^{-9}$, the pose
angles with the instrument's kinematics, and the two-point alignment against synthetic points.

Software development and the first draft of this text used an AI assistant (Claude, Anthropic) under the direction of
the authors, who reviewed the code and the text and are responsible for both. \tbd{confirm the scope of this statement
with all authors before submission}.

\section{Evaluation method}
\subsection{Platforms}
Microsoft HoloLens 2 (optical see-through), \tbd{operating system and browser versions}; Meta Quest \tbd{model},
\tbd{operating system and browser versions} (video see-through). Both load the same page.

\subsection{Technical measures}
Frame rate and 95th percentile frame time from the session log; alignment residual from the log; drift measured with a
ruler at the $20^\circ$ stop after one walk round the bench, five repeats; fraction of the light-arm arc in view from
0.75~m without moving the head; legibility of lines rated 1--5 with the laboratory lights on and off.

\subsection{User study}
\tbd{design, number of participants, ethics approval number}. Planned measures: learning gain between a pre-quiz and a
post-quiz, the System Usability Scale \parencite{X9}, the NASA Task Load Index \parencite{X10} and the Simulator
Sickness Questionnaire \parencite{X11}; the control condition is the same lessons in a desktop browser. No participant
data will be collected before approval.

\section{Results}
No measurement has been made. Table~\ref{tab:results} lists what will be reported.

\begin{table}[ht]
\caption{Comparison of the two platforms on identical content. All cells are empty until measured}\label{tab:results}
\begin{tabularx}{\linewidth}{Xlll}
\toprule
Measure & Unit & HoloLens 2 & Meta Quest\\
\midrule
AR session starts & yes / no & \tbd{} & \tbd{}\\
Frame rate & frames per second & \tbd{} & \tbd{}\\
Frame time, 95th percentile & ms & \tbd{} & \tbd{}\\
Alignment residual & mm & \tbd{} & \tbd{}\\
Drift after one walk round the bench & mm & \tbd{} & \tbd{}\\
Light-arm arc in view at 0.75 m & \% & \tbd{} & \tbd{}\\
Legibility & 1--5 & \tbd{} & \tbd{}\\
Learning gain & points & \tbd{} & \tbd{}\\
System Usability Scale & 0--100 & \tbd{} & \tbd{}\\
NASA Task Load Index & 0--100 & \tbd{} & \tbd{}\\
Simulator Sickness Questionnaire & score & \tbd{} & \tbd{}\\
\bottomrule
\end{tabularx}
\end{table}

\section{Discussion}
\tbd{write from the measured table}. Questions the data should answer: whether the narrower view of an optical
see-through display hides part of the light-arm arc at working distance; whether video passthrough makes thin overlays
easier to read under a bright lamp; whether fingertip pinches register the overlay well enough to read axis limits
against the machine's printed scale; and whether a web page reaches the frame rate of each device.

\section{Limitations}
The overlay is not occluded by the real machine. Alignment assumes a level bench. The lobe is a fitted model, not
measured data. The comparison uses a web runtime on both devices; a native build may perform differently.
\tbd{limitations found during measurement}.

\section{Conclusion}
\tbd{write after the results}.

\section*{Declarations}
\textbf{Funding} \tbd{funder and grant number, or none}.\quad
\textbf{Competing interests} \tbd{statement for every author}.\quad
\textbf{Ethics approval and consent} \tbd{board, approval number, consent procedure; required for the user study}.\quad
\textbf{Data availability} Source code, design files and session logs: \tbd{archive DOI}.\quad
\textbf{Author contributions} \tbd{roles, CRediT taxonomy}.

\printbibliography
\end{document}
